\section{Einleitung}
%TODO vlt noch irgendne wissenschaftliche Publikation zur Bedeutung von Datensicherheit bei KRITIS
Die zunehmende Digitalisierung im Gesundheitswesen bietet immense Chancen für effizientere Prozesse und die Verbesserung der Patientenversorgung.
In einem Umfeld, das als kritische Infrastruktur (KRITIS) eingestuft wird \cite{BSI_KRITIS}, steigen jedoch mit der zunehmenden Vernetzung auch die Anforderungen an die Sicherheit und Integrität der zugrunde liegenden Datenbankstrukturen weiter an. \\\\
Die steigende Zahl von Cyberangriffen auf Gesundheitseinrichtungen verdeutlicht, dass dringend robuste Sicherheitskonzepte nötig sind, da nicht mehr nur die Verfügbarkeit der Systeme, sondern gezielt die Vertraulichkeit hochsensibler Daten angegriffen wird. Ein aktuelles Beispiel hierfür ist der Vorfall bei den LUP-Kliniken in Mecklenburg-Vorpommern \cite{LUP_Leak}. Nach einem Ransomware-Angriff im Juli 2025 wurden sensible Patientendaten sowie Informationen über interne Abläufe im Darknet veröffentlicht. Der unkontrollierte Abfluss dieser Informationen ermöglicht beispielsweise gezielte Identitätsdiebstahle und Erpressungsversuche auf Basis sensibler Diagnosen.\\
Solche Vorfälle zeigen deutlich, dass der Schutz vor unbefugtem Datenzugriff keine rein theoretische Anforderung der DSGVO darstellt, sondern essenziell für das Vertrauensverhältnis zwischen Patienten und den Kliniken ist.\\\\
Das konzipierte System adressiert diese komplexen Anforderungen durch eine umfangreiche relationale Datenbankstruktur, die drei zentrale Bestandteile des Klinikalltags integriert. \\
Die Personalverwaltung bildet dabei die organisatorische Grundlage, indem sie die Differenzierung zwischen ärztlichem Dienst und Pflegepersonal sowie deren Zuordnung zu spezifischen Fachabteilungen sicherstellt.\\
Darauf aufbauend ermöglicht die Patienten- und Raumverwaltung eine präzise Steuerung der stationären Aufnahmen und eine effiziente Belegungsplanung.\\
Die dritte Säule bildet die Medikamenten- und Diagnosenverwaltung, welche die medizinische Historie der Patienten dokumentiert.\\\\
Aufgrund der hohen Sensibilität dieser Informationen liegt ein besonderer Fokus des Systems auf der Gewährleistung der Datensicherheit. Um Szenarien wie den oben genannten Datenleaks proaktiv zu begegnen, werden daher kritischer Datenfelder sicher verschlüsselt. Ziel ist es, ein System zu schaffen, das die Komplexität eines modernen Krankenhauses, wie beispielsweise der Uniklinik Würzburg, abbildet und gleichzeitig durch moderne Sicherheitsstandards die Vertraulichkeit der Patienten- und Personaldaten auf Datenbankebene schützt.